\newpage
\section{Clase de Investigación}
\noindent Este proyecto involucra investigación aplicada de tipo teórico-computacional, con el objetivo de explorar el fenómeno y producir nuevo conocimiento mediante el estudio y la simulación numérica de fenómenos físicos complejos, como la inestabilidad de Kelvin-Helmholtz en el ámbito de plasmas relativistas.

Además, dado que utilizamos herramientas de simulación para representar sistemas que no pueden reproducirse en el laboratorio, nos referimos a este estudio como uno de experimentación numérica. El cual pretende establecer y comprender las relaciones causales entre diversos parámetros físicos como son la resistividad el campo magnético ó la densidad entre otros y su influencia en el crecimiento o atenuación de las perturbaciones.

El enfoque adoptado es principalmente cuantitativo, ya que se trabajará con modelos matemáticos, códigos numéricos y análisis de datos derivados de computadoras, lo que permitirá una evaluación cuantitativa de cómo evoluciona el sistema de estudio y su comparación con la literatura del estado del arte, consignada en trabajos previos como los de \citep{mizuno-2013}, \citep{pimentel-2019}, \citep{la-rosa-2016}, \citep{tian-2016} y \citep{gomes-2017}.
